Este capítulo ha recorrido la construcción completa de SocratiCode, desde las decisiones de entorno y la elección de tecnologías hasta los detalles de implementación y la interfaz final. El resultado es un sistema en el que cada pieza cumple una función concreta dentro del objetivo pedagógico. El \textit{system prompt} define el comportamiento socrático del tutor sin necesidad de reentrenar ningún modelo. El flujo SSE permite que las respuestas lleguen al alumno de forma progresiva, manteniendo la sensación de conversación natural. El sistema de moderación dual protege el entorno educativo tanto frente a contenido inapropiado como frente a intentos de obtener soluciones directas. Piston aisla la ejecución del código del alumno en contenedores efímeros y cierra el ciclo al pasar la salida de la terminal como contexto al tutor. Por último, la interfaz SPA integra editor, terminal y chat en un único espacio visual que elimina la necesidad de alternar entre aplicaciones. En conjunto, estas decisiones de diseño e implementación materializan los requisitos funcionales definidos en el Capítulo~\ref{CAP:ANALISIS_DISENO} y sientan la base sobre la que se evaluará el sistema en el siguiente capítulo.
